% \thispagestyle{empty} % remove page number from this page

\begin{center}
	{\Large\bfseries ABSTRACT}
\end{center}

Automated crowd counting from visual data is a core capability for public-safety monitoring, transport-capacity management, and urban planning, yet the two main algorithmic approaches are each known for having a key drawback: detection-based counting can accurately count the number of people in a sparse-to-moderate scene, but fails when the scene is highly occluded, and density-regression counting can accurately estimate the number of people in a large, dense group but is less accurate in a small group that can be solved with reasonable accuracy. In this study, we propose, implement, and empirically test a sparse (detection-based) branch and a dense (density-regression) branch, both of which operate on the same input but are decidedly different in their objectives, jointly executed with a lightweight rule-based scene-density classifier and finally reconciled through a disagreement-gated fusion mechanism.  The full pipeline was then tested end-to-end on 285 images ranging from one to about 60 people in a ground-truth benchmark test set. The difficulty of setting up a detector at the head level by using the annotated positions significantly lowered the mean absolute error of the detector from 7.24 people per image to 3.36 people per image. The two density-regression networks, MDNN and CSRNet, achieved mean absolute errors of 12.94 and 10.66 respectively; both demonstrated the expected poor performance on sparse scenes. The fusion based on the CSRNet achieved more than half (MAE 4.34) the reduction in density-regression error, and for MDNN, more than fourfold (MAE 4.42) the reduction, compared to the single networks based fusion, and was much closer to the error obtained from the detector; on this basis the fusion based on the CSRNet was selected as the preferred hybrid, as the error in the worst case of the rare images where the detector failed, was significantly smaller. In neither hybrid configuration did the average-case accuracy of the fine-tuned detector outperformed when used alone, suggesting a significant operational advantage of a more resilient worst-case counting behaviour under failure for the hybrid detectors in this study, but not necessarily a corresponding improvement in average-case accuracy over the strongest individual detector; this finding does not diminish, but strengthens the argument for hybrid architectures in deployment environments where bounded worst-case error is of greater operational value than marginal average case gains.

\noindent\textbf{Keywords:} public transport; crowd counting; computer vision; object detection; density estimation; sensor fusion; scene classification; hybrid architecture;

\clearpage